% !TEX root = ../main.tex
% ============================================================
% 5. CONCLUSIONES Y TRABAJOS FUTUROS
% ============================================================
\chapter{Conclusiones y trabajos futuros}
\label{cap:conclusiones}

\section{Conclusiones}

\pendiente{una conclusión por objetivo específico, con el número que la
sostiene.}

% Repo: el texto que acordó el concilio del 21/10 (redacción A, reescrita).
%       Sus números: resultados/comparacion_cruzada/, resultados/cruzada_dlhm/,
%       resultados/barrido_obj2/ y resultados/hologramas_carlos/.
\paragraph{La pregunta de investigación.} El control explícito del muestreo
conviene cuando el campo propagado no cabe en la ventana del registro. Con $N$
píxeles de tamaño $\delta$, longitud de onda $\lambda$ y magnificación $M$, eso
ocurre cuando la distancia equivalente $z_\mathrm{eq}$ es del orden de
$z_\mathrm{lim} \approx N(\delta/M)^2/\lambda$ o mayor; con onda plana, $M = 1$
y $z_\mathrm{eq} = z$. En la DLHM, con $M = L/z$ y $z_\mathrm{eq} = z(L-z)/L$
(Ec.~\eqref{eq:escalado}), el cociente
$z_\mathrm{eq}/z_\mathrm{lim} = \lambda L(L-z)/(N\delta^2 z)$ crece al acercar
la muestra a la fuente: en los hologramas del laboratorio, a 0.4--2~mm de ella,
vale de 1.2 a 6.2 con el sensor entero y la geometría inferida, mientras que en
un DHM típico $N\delta^2/\lambda$ ronda los 46~mm y un reenfoque de milímetros
queda muy por debajo. La idea de que el MPASM rinde con mayores distancias se
sostiene, entonces, si la distancia se lee como equivalente, medida frente a
$z_\mathrm{lim}$, y no como física; y el criterio es un orden de magnitud, no
una frontera, porque sale del escalado paraxial. Por encima de él hay que
separar la propagación de la reconstrucción. Al propagar ---que es como se
simula un holograma--- el MPASM con sobremuestreo ($s = 3$) y $K_f$ automático
sigue a Rayleigh-Sommerfeld donde el FFT-ASM se contamina de aliasing, con
entre 10 y 1100 veces menos error; con $s = 1$ pierde frente al FFT-ASM con
relleno al doble. Al reconstruir, $K_f$ y $s$ no agregan información: con la
misma ventana, las dos cadenas dan casi la misma imagen de los hologramas
(correlación de 0.87 a 0.97), y subir $s$ de 1 a 4 cambia menos de 0.01~dB. Lo
que mejora la reconstrucción es una ventana mayor, y lo que aporta el MPASM es
poder pagarla: con onda plana, ampliarla le cuesta 32~MB frente a unos 32~GB al
FFT-ASM con relleno de ceros, y con los hologramas el MPASM esférico ---que
lleva la iluminación dentro de sus matrices y no usa $s$--- procesa en una GPU
de \SI{4}{GB} un recorte 20 veces mayor que el de la cadena del laboratorio,
con el que el contraste en los elementos de 5 a 10~\si{\micro\metre} casi se
duplica. Esa mejora la da la apertura, no el propagador, y sin referencia mide
contraste, no fidelidad.

\section{Trabajos futuros}

\pendiente{lo que quedó fuera del alcance o abierto por los resultados.}
